% Lösungen Einheit 2 von 2
\loesungskopf{le2}{Einheit 2 von 2 · Diagramme beurteilen und Boxplot}
\erg{24}{a) ja \quad b) nein, 80 \quad c) nein, $1{,}50$\,€ \quad d) ja}
\erg{25}{a) wahr \quad b) wahr (in Tausend) \quad c) falsch, 12\,000 \quad d) falsch, 41\,000 \quad e) falsch, von Juli (46\,000) zu August (41\,000) sank die Zahl \quad f) wahr, $46\,000 - 34\,000 = 12\,000$}
\erg{26}{a) wahr, $32\,\% + 20\,\% = 52\,\% > 50\,\%$ \quad b) wahr, $15\,\% + \frac{1}{3} \cdot 15\,\% = 20\,\%$ \quad c) falsch, jede 15. wären $\frac{1}{15} \approx 6{,}7\,\%$, nicht $15\,\%$ \quad d) falsch: der erste Teil stimmt ($24\,\% > 20\,\%$), der zweite nicht: $9\,\%$ von 600 sind 54 \quad e) wahr ($32\,\% + 9\,\% = 41\,\%$); falsch (jede 5. wären $20\,\%$, Sport hat $24\,\%$); nicht entscheidbar (die Umfrage trennt nicht nach Mädchen und Jungen)}
\erg{27}{a) 400 \quad b) Nein: von 420 auf 440 sind nur 20 Mitglieder mehr; die Säulen zeigen nur den Teil über 400 (20 und 40). \quad c) $20 : 420 \approx 4{,}8\,\%$ \quad d) Die Säulen werden doppelt so hoch, die Unterschiede wirken doppelt so groß, der Anstieg wirkt steiler; die Zahlen bleiben gleich. \quad e) Die Achse beginnt bei $2{,}30$\,€; die Säule der 5. Woche ist dreimal so hoch wie die der 1. Woche ($0{,}30$ gegen $0{,}10$ über dem Achsenanfang), der Preis stieg aber nur von $2{,}40$\,€ auf $2{,}60$\,€, also um $0{,}20 : 2{,}40 \approx 8{,}3\,\%$.}
\erg{28}{a) Die Zahl sinkt insgesamt von 120 auf 60 Zeitungen je Tag, also auf die Hälfte. \quad b) Nein: von 2021 (95) auf 2022 (100) stieg sie um 5. \quad c) Rückgang $120 - 60 = 60$ in 6 Jahren, also 10 je Jahr; $60 : 10 = 6$ Jahre nach 2025 wäre 2031. Unsicher, weil das Diagramm nur bis 2025 reicht und der Rückgang nicht gleichmäßig bleiben muss (2022 stieg die Zahl sogar).}
\erg{29}{a) 42; 40; 41; 37; 35 (Tausend) \quad b) Säulen mit diesen Höhen, Achse ab 0 \quad c) Links wirkt der Rückgang sehr groß, rechts sind die Säulen fast gleich hoch. \quad d) „jedes Jahr gesunken“ ist falsch: 2023 stieg die Zahl von 40 auf 41 Tausend. „fast halbiert“ ist falsch: $35 : 42 \approx 0{,}83$, ein Rückgang um etwa $16{,}7\,\%$. Die Achse beginnt bei 30; darum ist die Säule für 2025 (5 über dem Achsenanfang) weniger als halb so hoch wie die für 2021 (12).}
\erg{30}{a) 13 \quad b) 6 und 22 \quad c) 10 und 16 \quad d) $22 - 6 = 16$ \quad e) zwischen 12 und 15 \quad f) Die 7b hat den höheren Median (14 gegen 13). Bei der 7a liegen die Punkte weiter auseinander: Spannweite 16 gegen 10, Box 6 gegen 3 Punkte breit.}
\erg{31}{a) Box von 7 bis 12, Strich bei 10, Antennen bis 3 und 18 \quad b) Box von 6 bis 15, Strich bei 11, Antennen bis 5 und 20 \quad c) sortiert $2;\ 4;\ 5;\ 5;\ 6;\ 7;\ 8;\ 9;\ 9;\ 11;\ 13$: kleinster Wert 2, unteres Quartil 5, Median 7, oberes Quartil 9, größter Wert 13}
\erg{32}{a) Mo vergleicht nur die Säulenhöhen; die Achse beginnt bei 200, die Säulen zeigen 60 und 120. Richtig: $320 - 260 = 60$ Besucher mehr, $60 : 260 \approx 23\,\%$ mehr, nicht doppelt so viele. \quad b) wirkt 3-mal so hoch (15 gegen 5 über dem Achsenanfang 50); wirklich $10 : 55 \approx 18{,}2\,\%$ mehr}
\erg{33}{a) Die Höhe einer Säule soll die ganze Anzahl zeigen; nur dann heißt doppelt so hoch auch doppelt so viel. Beginnt die Achse später, werden kleine Unterschiede übertrieben. \quad b) Nein: Hier kommt es auf kleine Änderungen an, und eine Temperatur von $0\,°$C kommt bei einem Menschen nicht vor; die Werte sind Punkte an einer Linie, keine Säulen, deren Höhen man vergleicht. Die Achse muss dann aber deutlich beschriftet sein.}
\erg{34}{a) $125 - 5 \cdot 20 = 25$ Plätze fehlen; 2 Ständer (1 Ständer reicht nur für 20) \quad b) $(125 - 84) : 84 \approx 48{,}8\,\%$ \quad c) Nicht nötig und nicht fair: Die Zahl ist um fast die Hälfte gestiegen, das sieht man auch mit einer Achse ab 0. Eine Achse ab 80 würde den Anstieg übertreiben und den Antrag unglaubwürdig machen.}

\noindent Lösungsgrafik 31:

\begin{boxplots}[xmin=0,xmax=20,xstep=1,xlabel=Wert]
\bp{2,5,7,9,13}{c)}
\bp{5,6,11,15,20}{b)}
\bp{3,7,10,12,18}{a)}
\end{boxplots}
